% !TEX root = ../main.tex
\chapter{Models: train on what Analyse produced}\label{ch:models}
Models answers one question honestly: can a classifier learn what you mean by an interesting
\term{window}{window}, and do labels made by clustering teach it as well as labels made by a person?
It trains on windows Analyse cut, scores the result on windows it never saw, and sets every score
against a \term{null}{null}.

\begin{plainterms}[An exam the student has not seen]
A classifier is like a student. Marking it on the questions it practised says nothing; you mark it on
an exam it has never seen. Models keeps three such exams apart: (i) the last stretch of time on each
channel it trained on, (ii) channels it never trained on, and (iii) a whole recording held back until
the end of the project. And it compares the student with one who was given shuffled answers: if both
score the same, the first one learned nothing.
\end{plainterms}

The workspace has four tabs that follow one job from start to decision. \apppage{models.launch} sets
up and starts a \emph{paired} training job; \apppage{models.results} reads one arm of it on one exam;
\apppage{models.compare} sets the two arms against each other; \apppage{models.registry} is where a
model is signed off for use. Launch, Results and Compare read the bridge and the core's paired training
job (\code{Working/training/}); the Registry is still a preview on example data. A \emph{Jobs} link on
every tab opens \cref{ch:jobs}.

\begin{plainterms}[Paired label arms]
A paired job trains the same classifier twice, on exactly the same windows, with the same settings and
the same random seed. Only the labels differ. \textbf{Arm A} learns from your manual verdicts (the
labels a person gave in Review). \textbf{Arm B} learns from cluster labels: the training windows are
grouped by a clustering, and each group is translated to \emph{interesting} or \emph{not
interesting}. Because nothing else differs, any difference in their scores can be put down to the
label source. The classifier in both arms is a random forest, which is also the baseline.
\end{plainterms}

% ---------------------------------------------------------------------------------------------
\pagesection{models.launch}{Launch}{Train a template across channels, with paired label arms}{live}
Launch sets up one paired training job in five numbered cards, checks it, estimates its cost, and
either trains it on this machine or writes a SLURM script for the cluster.

\pagefield{When you'd come here}
\begin{itemize}
  \item You saved a training template in Analyse (a chain with a window-matrix step) and want to
        train on several channels.
  \item You pressed \ui{Train model} on a Window pool chain in Analyse: Launch opens with arm B.2
        filled in (see below).
\end{itemize}

\screenshot{pages/models/models.launch--live-default}{The five cards of a paired job: template,
sources, label arms, evaluation and options, with the checks and the cost estimate.}

\pagefield{Doing things here}
\begin{enumerate}
  \item \textbf{1 · Training template.} Pick a template. Its window-matrix step sets the window
        length, the step and the features. \ui{Open in Analyse} opens it on the chain page.
  \item \textbf{2 · Sources · channels.} Give each channel of the recording a role: \ui{train}
        (cut into time blocks; its last blocks are exam i), \ui{exam} (scored, never trained on:
        exam ii) or \ui{off}. Set the gap between blocks (\ui{\(\ge\) 1 window} or
        \ui{\(\ge\) 2 windows}). Press \ui{Save window set} to save the pooled windows with their
        blocked split, or pick a saved set from the table. The held-out recording is shown locked.
  \item \textbf{3 · Label arms.} Arm A needs nothing. For arm B press \ui{Propose cuts}: the training
        windows (test windows take no part) are clustered once, and a table lists each number of
        clusters \(k\) with its silhouette score, cluster sizes and how many clusters are real rather
        than specks. Pick a \(k\) (one is marked \texttt{draft}). A translation table then maps each
        cluster to \emph{interesting} or \emph{not interesting}, drafted by the majority human verdict
        on the training windows, with each cluster's purity; a cluster under 75\,\% is flagged
        \texttt{impure}. Change the translation now.
  \item \textbf{4 · Evaluation.} Shows the split by time for each channel, and the three exams,
        never pooled.
  \item \textbf{5 · Options.} The forest's size and seed (shared by both arms), the number of
        label shuffles for the null, and the bootstrap.
  \item \ui{Before launch} lists the checks the job itself runs before it starts; an error disables
        both buttons and names itself. The estimate chip gives the cost from the measured time of one
        forest fit times the number of fits.
  \item Press \ui{Train locally} when the estimate is within the local limit, or
        \ui{Create SLURM script} when it is over. When a local job finishes, \ui{Open Results}.
\end{enumerate}

\begin{watchout}[The cut freezes after the first test score]
Once any run has scored a window set's test block at some \(k\), the cut and the translation table
are frozen on that set, and another \(k\) is refused. Choosing \(k\) after seeing test scores would
be tuning on the exam. Decide the cut and the translation before you train.
\end{watchout}

\subsection*{Arm B.2: from a Window pool}
\ui{Train model} on an Analyse Window pool chain ending in Shape clustering opens Launch as
\emph{arm B.2 · cluster labels · trace shape}. The template and the pool arrive filled in; the cut and
the mapping are read-only here because they are set on the Shape clustering block's page in Analyse,
where the dendrogram is. The pool is unlabelled, so there is no arm A. \ui{Train locally} trains;
\ui{Create SLURM script} writes the cluster job instead.

\pagefield{Where it leads}
\apppage{models.results} for the scores; Jobs for a cluster job; \apppage{analyse.chain} via
\ui{Open in Analyse}.

\begin{underthehood}
Client: \code{webui/client/src/models/LaunchPage.tsx} (\code{PairedLaunchPage}),
\code{B2Launch.tsx}; reads in \code{webui/client/src/api/models.ts} (\code{getModelsSetup},
\code{proposeCut}, \code{checkPairedJob}, \code{trainPairedJob}, \code{slurmPairedJob}). Bridge:
\code{webui/server/training_routes.py} (\code{/api/models/setup}, \code{/windowsets},
\code{/propose}, \code{/checks}, \code{/train}, \code{/slurm}); arm B.2 in
\code{webui/server/shape_routes.py} (\code{/api/models/b2/setup}, \code{/b2/train}). Arm A's labels
are \blockref{catalogue.manual_labels}; the classifier is \blockref{catalogue.classifier}.
\end{underthehood}

% ---------------------------------------------------------------------------------------------
\pagesection{models.results}{Results}{One arm on the test block, against baseline and nulls}{live}
Results shows one paired run, one arm and one exam at a time, never pooled. Each number comes with
how sure it is and with what chance would give.

\pagefield{When you'd come here}
When a training job has finished, to see how well each arm did, and whether it beat its null.

\screenshot{pages/zz_models_ab/models.results--fixup-ab-arm-a}{Arm A on exam i: macro F1 with its
interval, the null behind it, the confusion matrix, per-class and per-channel scores, and
calibration.}

\pagefield{Doing things here}
\begin{enumerate}
  \item Pick the run, the arm (A or B) and the exam (i, ii, or iii; exam iii stays locked).
  \item Read the headline tiles: \emph{macro F1} (the average of the F1 score of each class, so a
        rare class counts as much as a common one) with its block-bootstrap confidence interval, and
        balanced accuracy.
  \item \ui{Against the null}: the arm's macro F1 (blue, with its band) over the spread of macro
        F1s of the same forest trained on shuffled labels. An arm inside that spread learned nothing
        usable.
  \item \ui{Confusion}: rows are the human verdict, columns the arm's prediction, row-normalised,
        with counts. \ui{Per class} gives precision, recall and F1 with \(n\); \ui{Per channel} gives
        both arms' macro F1 on each channel.
  \item \ui{Calibration · validation block}: how well the model's \(P(\text{interesting})\) matches
        reality on the validation block, and the lowest threshold that reaches a target precision
        there.
  \item \ui{Reference line · the existing MODELS/} scores older models on the same exams. They were
        trained differently, probably on overlapping labels, so they are a reference, not an arm.
  \item \ui{Yardstick (B)} scores arm B in its own cluster vocabulary against blind hand-labelling of
        test windows.
\end{enumerate}

\begin{plainterms}[Calibration]
A weather forecast that says ``70\,\% chance of rain'' is calibrated if it rains on about 70\,\% of
those days. A calibrated model's ``70\,\% interesting'' is right about 70\,\% of the time. The
reliability plot shows how close the model comes.
\end{plainterms}

For an arm B.2 run, Results shows the blind test instead: \ui{Label test windows blind} opens a Review
queue of the run's test windows with no hint of the model's call; afterwards the page sets your answers
against the model's, per exam, beside how often you agree with yourself on windows shown twice (no
model can be expected to beat that), and beside the manual-label models scored on the same windows.

\pagefield{Where it leads}
\apppage{models.compare} for the paired difference; \apppage{models.registry} to register a model.

\begin{underthehood}
\code{webui/client/src/models/ResultsPage.tsx}, \code{B2Blind.tsx}; \code{GET /api/models/runs} and
\code{/api/models/runs/{run_id}} (\code{webui/server/training_routes.py}); the blind test in
\code{webui/server/blind_routes.py}.
\end{underthehood}

% ---------------------------------------------------------------------------------------------
\pagesection{models.compare}{Compare}{Manual against cluster labels, paired}{live}
Compare sets arm A against arm B of one paired run. The only thing that differs between them is the
label source, so the difference is attributable to it.

\pagefield{When you'd come here}
To answer whether cluster-derived labels teach a classifier as well as manual ones (research question
one), with a measure of how sure the answer is.

\screenshot{pages/zz_models_ab/models.compare--fixup-ab-a-vs-b}{Arm A against arm B: the paired
difference in macro F1 with its interval, McNemar's test, and the window-by-window agreement.}

\pagefield{Doing things here}
\begin{enumerate}
  \item Pick the run and the exam.
  \item \ui{Paired difference}: \(\Delta\)F1 between the arms under the same bootstrap resamples,
        with its confidence interval; the line at zero is no difference. McNemar's test asks whether
        the two arms are wrong on different windows more often than chance.
  \item \ui{Agreement, window by window}: a 2\(\times\)2 table of \emph{both right}, \emph{only A
        right}, \emph{only B right}, \emph{both wrong}. Click a cell to step through its windows with
        \ui{Prev} and \ui{Next}; \ui{Open the channel in Explore} shows a window in context.
  \item Per class and per channel differences.
  \item \ui{Cluster \textrightarrow{} class}: how arm B's clusters fall against the human verdicts on
        the training windows, the translation table fixed in the recipe, and purity.
\end{enumerate}

\begin{watchout}[An interval that crosses zero]
If the \(\Delta\)F1 interval includes zero, the data do not tell the two label sources apart on that
exam. That is a result, not a failure.
\end{watchout}

\pagefield{Where it leads}
Back to \apppage{models.results}; \apppage{explore.signal} for a window in context.

\begin{underthehood}
\code{webui/client/src/models/ComparePage.tsx}; \code{GET /api/models/runs/{run_id}/compare} and
\code{/disagreements} in \code{webui/server/training_routes.py}.
\end{underthehood}

% ---------------------------------------------------------------------------------------------
\pagesection{models.registry}{Registry}{Register with held-out checks and sign-off}{preview}
The registry lists every model the installation has scored, with its status: \emph{candidate}
(results, no sign-off), \emph{registered} (can be inserted in Analyse as a Model stage),
\emph{retired} (can be restored) or \emph{rejected} (keeps the checks it failed). Registering a
candidate is a three-part gate: automatic held-out checks, a sample of its windows judged blind in
Review, and a written decision.

\textbf{This page is a preview.} Its models come from \code{webui/client/src/fixtures/models.ts};
register, retire, restore and reject change only browser memory, and \ui{Open sample in Review} makes
a pretend queue. No bridge route serves a model registry yet.

\pagefield{When you'd come here}
To see how a model is meant to be signed off before Analyse may use it.

\screenshot{pages/models/models.registry--candidate}{A candidate's registration card: its held-out
checks, the blind verification sample, and the sign-off form.}

\pagefield{Doing things here}
\begin{enumerate}
  \item Filter the list: \ui{all}, \ui{registered}, \ui{candidates}, \ui{retired}. Each row says what
        uses the model.
  \item Open a candidate. \textbf{1 · Held-out checks} are automatic: a failure blocks registration,
        a warning needs a written reason.
  \item \textbf{2 · Verification}: judge a sample of the model's windows in Review (\ui{Open sample in
        Review}, \ui{Add 20 more}). Sign-off waits until the sample is judged.
  \item \textbf{3 · Decision}: give the registered name and version and your reason, then register, or
        \ui{Reject}. Only this third step registers anything.
  \item A registered model can be retired (blocked while a template uses it) and later restored with
        its old sign-off.
\end{enumerate}

\pagefield{Where it leads}
A registered model would be offered in Analyse's \ui{Insert a stage} window as a Model stage, its
calibration thresholds becoming the threshold stage's recommended values \unverified{(this is the
design; no live code path does it yet)}.

\begin{underthehood}
\code{webui/client/src/models/RegistryPage.tsx}; \code{getRegistry} in
\code{webui/client/src/api/models.ts} returns fixture data through \code{demo(...)}.
\end{underthehood}
